\par\Needspace{10\baselineskip}\pdfbookmark[1]{Pasquich. 1817 October 4.}{pasquich1817}
\begin{center}
\rule{0.65\linewidth}{0.4pt}\\
Göttingische gelehrte Anzeigen. 1817 October 4.\\
\rule{0.65\linewidth}{0.4pt}
\end{center}

\textit{Abgekürzte Logarithmisch-Trigonometrische Tafeln, mit neuen Zusätzen zur
Abkürzung und Erleichterung trigonometrischer Rechnungen, herausgegeben von }\textsc{Joh.}
\textsc{Pasquich}\textit{, Director der Königl. Ofner Sternwarte.} Leipzig. At the Weidmann bookshop, 1817. XXXII and 228 pages in octavo. (Also with a Latin title.)
\par\SourceLeaf{256}
Smaller logarithm tables with five decimal places are very popular among those who do much numerical calculation, especially astronomers. In fact, the cases in which they suffice are frequent, indeed the more frequent ones, and a convenient format and moderate size make the work much easier. \textsc{Lalande}'s neat small tables and \textsc{Prasse}'s tables are in everyone's hands. In the latter, compression into a small space has been carried as far as possible, partly, however, at the expense of convenience. As the preface reports, the publication of the present tables, which also give only five places, was prompted by the wish expressed by \textsc{Gauss} in the Monatliche Correspondenz of 1812: that the table first printed there for conveniently calculating the logarithms of sums and differences might be incorporated into a new edition of \textsc{Lalande}'s tables. The reprinting of this auxiliary table in the present collection will be welcome to those who had no access to the first printing, or found its use in the M. C. too troublesome. This collection is also distinguished by a new auxiliary table calculated by Mr. \textsc{Pasquich} and appended to the trigonometric tables, which we shall discuss in more detail below.

As in \textsc{Lalande} and von \textsc{Prasse}, the logarithms of numbers extend to 10000 and are printed here, as there, in their natural order. Two conveniences found in \textsc{Lalande}'s tables are unfortunately missing: the \textit{differences} are not supplied, and the lowest logarithms in each column are not repeated at the top of the next. These conveniences are absent from \textsc{Prasse}'s tables too, but there their loss is more than compensated by the small space occupied by the table: 24 pages there give exactly what takes 56 pages of a considerably larger format in \textsc{Pasquich}. These may seem trifles, and they are for anyone who only looks up logarithms now and then, but not for those who must constantly have logarithm tables to hand.

\Needspace{9\baselineskip}In the trigonometric tables, each page on the left contains the logarithms of the sines, cosines, tangents and cotangents, with three pages covering two degrees. This arrangement, dictated by the chosen format and type, seems somewhat inconvenient to us; \SourceLeaf{257}
we would have preferred either a smaller format, with half a degree on each page, or one slightly longer, with smaller type, so that a whole degree would fit on a page, as in \textsc{Shervin}'s tables. Of these logarithms, only the last four, three or two digits are printed for as long as the preceding digits remain unchanged. This made the work easier for the copyist, typesetter and proofreader, and gives the tables a cleaner appearance. Nevertheless, we cannot unreservedly approve this arrangement for tables intended for daily use: the eye must always make the effort, however small, of going up the column to find the remaining digits. The differences of the logarithms are given here already divided by 60, an arrangement also chosen in some other tables with the intention of making interpolation easier. Whether it really does so will depend on the calculator's habits. In this respect the reviewer finds \textsc{Lalande}'s tables, which give the full differences, at least no less convenient. Since the numbers involved are small, a calculator with some practice can easily carry out the operations needed for interpolation mentally, and can almost always find some advantage peculiar to the case. There is also the pleasant certainty of obtaining the interpolation result as accurately as possible. With the arrangement chosen by Mr. \textsc{Pasquich}, on the other hand, the interpolation error is, generally speaking, somewhat greater; this is not the place to develop that point in detail.

\Needspace{24\baselineskip}In the trigonometric tables, the page on the right contains the squares of the sines, cosines, tangents and cotangents. These are intended to make interpolation easier when, from the logarithm of a sine, cosine, tangent or cotangent, the logarithm of one of the other three trigonometric functions is required without needing the arc itself. This operation does indeed occur extremely often. The ordinary method requires one multiplication and one division in interpolating, whereas with Mr. \textsc{Pasquich}'s auxiliary table one multiplication suffices. The following is sufficiently accurate for interpolation:
\SourceLeaf{258}
\[
\begin{aligned}
\Delta\log\cos\varphi&=-\operatorname{tang}\varphi^{2}.\Delta\log\sin\varphi\\
\Delta\log\operatorname{tang}\varphi&=-\Delta\log\operatorname{cotang}\varphi=(1+\operatorname{tang}\varphi^{2}).\Delta\log\sin\varphi\\
\Delta\log\sin\varphi&=-\operatorname{cotang}\varphi^{2}.\Delta\log\cos\varphi\\
\Delta\log\operatorname{tang}\varphi&=-\Delta\log\operatorname{cotang}\varphi=-(1+\operatorname{cotang}\varphi^{2}).\Delta\log\cos\varphi\\
\Delta\log\sin\varphi&=\cos\varphi^{2}.\Delta\log\operatorname{tang}\varphi=-\cos\varphi^{2}.\Delta\log\operatorname{cotang}\varphi\\
\Delta\log\cos\varphi&=-\sin\varphi^{2}.\Delta\log\operatorname{tang}\varphi=\sin\varphi^{2}.\Delta\log\operatorname{cotang}\varphi
\end{aligned}
\]
The reviewer must nevertheless confess that he finds the ordinary method of interpolation not merely as convenient, but even more convenient. For one thing, it always takes some effort to become so accustomed to the above six formulas that they can be applied correctly without any thought or without keeping a separate sheet beside one. For another, it is troublesome first to look up the multiplier on the opposite page, or rather to piece it together, since the separation of initial and final digits mentioned above has also been used in printing this table. Finally, the ordinary method always involves only small numbers, which are easy to calculate with mentally, whereas the squares in \textsc{Pasquich}'s tables are given to five decimal places. Not all these are needed, of course, but for that very reason they are distracting, as every experienced calculator knows. Moreover, we cannot leave unmentioned that the ordinary method is, generally speaking, \textit{more accurate} than this more elaborate interpolation. The reasons for this assertion, of which a less careful examination might perhaps suggest the exact opposite, would take too much space here. We shall content ourselves with illustrating what has been said by an example. If, for \(\log\cos\varphi=9{,}92478\), \(\log\operatorname{tang}\varphi\) is required, the proportional part is found from \textsc{Pasquich}'s table by calculating \(4\times(1+2{,}4170)=13{,}668\), or, rounded, \(=14\), hence \(\log\operatorname{tang}\varphi=9{,}80850\). The ordinary method gives the proportional part just as conveniently by evaluating \(\frac{4\times28}{9}=12\frac49\), rounded, \(=12\), and the required logarithm \(=9{,}80848\). In this example the ordinary method also gives the more accurate result. In other cases the reverse can hold, but on average the advantage in this respect will lie with the ordinary method. We do not wish to deny, however, that this part of the table, even if it does not benefit interpolation, may sometimes be welcome for \SourceLeaf{259}
other purposes. But the convenience of logarithmic hand tables intended for daily use naturally decreases in the same proportion as their bulk increases. We also note that, in the first degree, the trigonometric logarithms are given every \(10\) seconds in successive \(10\)-second intervals up to \(56\) minutes, and in the last four minutes every \(20\) seconds in successive \(20\)-second intervals.

\textsc{Gauss}'s table for the logarithms of sums and differences is reprinted entirely unchanged. It seems inconsistent to us, however, that the author questions in his introduction the usefulness of a similar table with seven decimal places. Provided such a table is suitably arranged, its usefulness in accurate calculations is just as great as that of the tables reprinted here in calculations with five decimal places. With the smaller tables, as with the larger, the saving in time will naturally be felt only by people who have \textit{much} calculation to do. We may now soon expect the appearance of such a larger table, calculated by a skilful hand.

